%%%% CAPÍTULO 3 - METODOLOGIA  (PROPOSTA para revisão)
%% Observações:
%%  - As remissões "(seção 2.3)", "(seção 3.2)" etc. estão em texto literal (provisório);
%%    serão convertidas em \ref{} quando o Capítulo 2 estiver escrito e com rótulos.
%%  - Marcadores de FIGURA e de CÓDIGO estão em comentário, para inserirmos depois.

\chapter{Metodologia}\label{cap:metodologia}

\section{Caracterização da pesquisa}

Quanto à natureza, este trabalho caracteriza-se como uma pesquisa aplicada, pois tem por finalidade a geração de um produto de uso prático (um material didático-multimídia voltado ao ensino de demonstrações em Geometria Plana), apoiando-se, para isso, em pesquisa bibliográfica. Quanto à abordagem, é uma pesquisa qualitativa, uma vez que se volta à concepção, à descrição e à análise interpretativa do produto desenvolvido, e não à mensuração de variáveis \cite{gerhardt2009metodos,ludke1986pesquisa}.

Quanto aos objetivos, a pesquisa é exploratório-descritiva: exploratória por buscar na literatura a fundamentação para o problema e para as soluções adotadas, e descritiva por detalhar sistematicamente a estrutura, o conteúdo e o processo de desenvolvimento do produto. Quanto aos procedimentos, apoia-se na pesquisa bibliográfica e no desenvolvimento de um produto educacional, combinando o levantamento das referências teóricas (Capítulo 2) com o relato documentado do processo de concepção e construção do material (seção 3.2).

O foco do trabalho está na concepção e na fundamentação do produto. Sua avaliação, descrita na seção 3.3, é de natureza teórica: examina a coerência entre o design do material e os referenciais que o fundamentam.

\section{Desenvolvimento do produto}

O desenvolvimento do produto seguiu quatro etapas, descritas a seguir: a seleção e a organização dos conteúdos, a produção do material escrito, a produção das animações e a integração entre texto e vídeo na plataforma.

\subsection{Seleção e organização dos conteúdos}

A definição dos conteúdos que comporiam o módulo inicial da trilha tomou como base duas obras de referência: \citeonline{dolcepompeo2013} e \citeonline{caminha2022geometria}. A opção por essas fontes justifica-se por consolidarem de modo rigoroso os fundamentos da Geometria Plana e por partirem de um mesmo ponto inicial, o que permitiu organizar o material sobre uma base amplamente reconhecida e já consagrada no ensino.

Optou-se por abordar, nesta etapa, apenas a porção inicial desse conteúdo, a saber, ângulos, triângulos e seus critérios de congruência, quadriláteros e a circunferência, de modo a fundamentar solidamente o começo da trilha. Trata-se de uma delimitação deliberada, coerente com o caráter de módulo inicial do produto: pretende-se, em trabalhos futuros, ampliar a trilha até abranger a totalidade da Geometria Plana.

A ordenação dos temas preservou a sequência adotada nas obras de referência, orientada pela construção lógica da Geometria: cada resultado é demonstrado utilizando apenas fatos previamente estabelecidos ou propriedades elementares assumidas como ponto de partida. Com isso, o estudante percorre um caminho em que nenhuma demonstração depende de conhecimentos ainda não apresentados, organização que favorece a progressão do raciocínio, do reconhecimento intuitivo à dedução formal (seção 2.3).

\subsection{Produção do material escrito}

As demonstrações foram adaptadas das obras de referência adotadas (seção 3.2.1), com reescrita orientada por dois objetivos simultâneos: preservar o rigor e a coerência lógica das provas e, ao mesmo tempo, torná-las adequadas à representação animada. Já na elaboração do texto, cada demonstração foi pensada em função dos elementos que poderiam ser visualmente destacados, como triângulos que se deslocam, retas que mudam de posição e construções auxiliares que surgem passo a passo, de modo que texto e animação fossem concebidos de maneira integrada, e não como peças independentes.

Essa integração desde a concepção é uma característica central do produto: a animação não ilustra um texto pronto, mas faz parte do mesmo projeto didático. Em todo o processo, a preocupação foi manter a clareza para o estudante que lê de forma autônoma, explicitando passos que as obras de referência frequentemente deixam implícitos.

O material escrito foi produzido em HTML, com as expressões matemáticas redigidas em \LaTeX{} e renderizadas pela biblioteca MathJax, o que garante boa legibilidade das fórmulas sem prejudicar a leitura. Cada tópico segue uma estrutura fixa: enuncia-se o teorema ou a definição e, em seguida, apresenta-se a demonstração em passos, acompanhada das figuras correspondentes e alinhada à sequência das animações.

\subsection{Produção das animações}

As animações foram produzidas na plataforma Conexty\footnote{Disponível em: \url{https://conexty.com}. Acesso em: 16 out. 2026.}, um ambiente online de criação de videoaulas no qual o autor programa, em JavaScript, os elementos geométricos a serem desenhados e animados sobre uma tela, combinando-os a texto e áudio e exportando o resultado em vídeo. A Figura~\ref{fig:conexty} apresenta esse ambiente, com o código, à esquerda, e a tela em que a animação é construída, à direita. Para cada demonstração, a animação foi codificada diretamente na plataforma, com a inserção de pausas que segmentam o raciocínio em etapas controláveis.

%% ============================================================
%% FIGURA A INSERIR (3.2.3) -- rotulo \label{fig:conexty}
%%   captura do ambiente do Conexty (codigo em JavaScript a esquerda e
%%   canvas/animacao a direita). JA CITADA no texto acima.
%% ============================================================

A construção seguiu um princípio didático constante: primeiro \emph{justifica-se} o passo (enunciando, em texto na tela ou na narração, o que será feito e por quê) para só então \emph{realizá-lo} visualmente. Assim, a marcação de um ângulo, o traçado de um segmento ou o destaque de um triângulo são sempre precedidos de sua justificativa, reforçando a compreensão do raciocínio, e não apenas sua visualização. O texto escrito na tela foi empregado de forma seletiva: nas deduções de maior impacto, a justificativa aparece simultaneamente escrita e narrada; nas passagens mais simples, mantém-se apenas a narração, para não sobrecarregar a tela. O Quadro~\ref{qua:codigo} traz um excerto de uma dessas animações, no qual esse padrão de justificar antes de marcar pode ser observado no próprio código.

\begin{lstlisting}[language=JavaScript,caption={Excerto de código de uma animação no Conexty (arco capaz): cada passo é precedido de um comentário que o descreve, e \texttt{pause()} segmenta as etapas.},label={qua:codigo}]
// por A, a semirreta que forma 90-alpha com AB (ate O); angulo 90-alpha em A
drawSegment({ points: [A2, O2], width: 0.024, color: cor_azul_eletrico })
drawSector({ points: [B2, A2, O2], radius: 0.5, fill: true,
            fillColor: cor_azul_eletrico, opacity: 0.5 })
drawText({ text: "$90^\\circ-\\alpha$", x: A2.x + 0.45, y: A2.y - 0.5,
          color: cor_azul_eletrico })

pause()

// centro O (encontro com a mediatriz)
drawPoint({ x: O2.x, y: O2.y, size: 0.08, color: cor_rosa_choque })
drawText({ text: "$O$", x: O2.x + 0.2, y: O2.y + 0.05, color: cor_rosa_choque })
\end{lstlisting}

O código completo dessa animação, com todas as etapas da construção, encontra-se no Apêndice~\ref{ap:arcocapaz}.
%% OBS: o pacote `listings` precisa ser carregado em configuracoes.tex ao aplicar no oficial.

Concluída a animação, o vídeo final foi gerado executando-a na própria plataforma com o acréscimo da narração do autor, que explica oralmente cada etapa, e exportado em formato de vídeo. As figuras estáticas que acompanham o texto escrito foram obtidas a partir das mesmas animações, garantindo coerência visual entre o que o estudante lê e o que assiste; a única diferença é o fundo (escuro nos vídeos e claro nas figuras do texto), o que exigiu pequenos ajustes de cor para preservar a legibilidade.

Quanto ao formato, optou-se deliberadamente por vídeos curtos, de até cerca de 180 segundos, em orientação vertical. Essa decisão alinha-se ao objetivo de engajamento: trata-se do formato predominante no consumo em dispositivos móveis e redes sociais, o que favorece o alcance a um público amplo (seção 2.5).

\subsection{Integração do texto e do vídeo na plataforma}

A trilha foi disponibilizada em uma plataforma digital de acesso público, o sítio \emph{Geometria Demonstrada}, em que cada tópico reúne, em uma mesma página, o vídeo e o texto correspondentes: o vídeo é apresentado à esquerda e a demonstração escrita, à direita, lado a lado. Essa disposição materializa o princípio da contiguidade da Teoria Cognitiva da Aprendizagem Multimídia (seção 2.4), segundo o qual palavras e imagens relacionadas devem ser apresentadas próximas, reduzindo o esforço de o estudante alternar entre fontes distintas de informação. A Figura~\ref{fig:pagina} ilustra essa disposição em uma das páginas do sítio, com o vídeo à esquerda e o texto da demonstração à direita.

%% ============================================================
%% FIGURA A INSERIR (3.2.4) -- rotulo \label{fig:pagina}
%%   captura da página do sítio "Geometria Demonstrada" com o
%%   vídeo à esquerda e o texto à direita. JA CITADA no texto acima.
%% ============================================================

Do ponto de vista técnico, os vídeos foram hospedados em uma lista de reprodução no YouTube e incorporados, por meio de link, à página correspondente da plataforma. Essa escolha facilita o acesso e a distribuição do material e mantém texto e vídeo reunidos no mesmo ambiente, à disposição de qualquer pessoa, a qualquer momento.

\section{Procedimentos de análise}

Como este trabalho não prevê aplicação com estudantes, a avaliação do produto é de natureza teórica e interpretativa: consiste em confrontar as decisões de design do material com os referenciais que o fundamentam (Capítulo 2), discutindo em que medida o produto é coerente com o que a teoria recomenda.

O procedimento adotado foi o seguinte. Primeiro, selecionaram-se instâncias representativas do produto, em especial demonstrações cujas animações concentram os recursos mais característicos da proposta, como o Teorema de Pitot e o arco capaz. Em seguida, cada instância foi examinada à luz de cada referencial: verificou-se se o sequenciamento dos conteúdos respeita a progressão do pensamento geométrico (Van Hiele); se a condução das demonstrações reflete as etapas da resolução de problemas (Pólya); se a combinação de texto e vídeo e a segmentação em passos curtos estão de acordo com os princípios da aprendizagem multimídia (Mayer); e se o formato dos vídeos atende ao propósito de divulgação e engajamento. Por fim, discutiram-se tanto as escolhas que se mostram bem fundamentadas quanto os limites e as tensões entre os princípios.

Trata-se de uma análise qualitativa, na qual o pesquisador atua como principal instrumento de investigação \cite{ludke1986pesquisa}, interpretando o próprio produto à luz da teoria. Os resultados dessa análise são apresentados no Capítulo 4.
